% ch03  §4.3–§4.7 / §5.1（原文 1296–1647 行）

\subsection{分解结构}\label{ss:geom fact}

对一簇 $U^{I}$ 上的几何对象 $A_{I}$（叠或层），每个有限集 $I$ 对应其中一个。所谓 $\{A_{I}\}$ 上的分解结构，通常包含两个方面。其一是 $A_{I}$ 限制到 $U^{I}_{\disj}$ 时的行为：例如当 $A_{I}$ 是 $U^{I}$ 上的叠时，可要求 $A_{I}|_{U^{I}_{\disj}}$ 配有与 $(\prod_{i\in I}A_{\{i\}})|_{U^{I}_{\disj}}$ 的同构。其二是 $A_{I}$ 在部分对角上的行为：例如可要求 $A_{I}|_{\D(U)}\cong A_{\{1\}}$。这类分解结构的典型例子由 $X^{I}$ 上的 Beilinson--Drinfeld Grassmannian $\Gr^{I}$ 承载。

在 shtuka 的情形，上述第一个方面只在弱的意义下成立，其内容本质上就是命题 \ref{p:fm sm} 与命题 \ref{p:local model}。下面把第二个方面写清楚。

对任意叠 $\frX$ 与有限集之间的任意映射 $\th: I\to J$，有诱导映射
\begin{equation*}
\D_{\th, \frX}: \frX^{J}\to \frX^{I}
\end{equation*}
它把 $(x_{j})_{j\in J}$ 送到 $(x_{\th(i)})_{i\in I}$。当 $\th$ 是满射时，$\D_{\th,\frX}$ 是 $\frX^{I}$ 中的部分对角。

Hecke 叠的如下性质由定义直接看出。若 $\th: I\surj J$ 是有限集之间的满射，则存在典范同构
\begin{equation}\label{Hk diagonal}
\Hk^{I}_{G,\bK}|_{\D_{\th,U}(U^{J})}\cong \Hk^{J}_{G,\bK}.
\end{equation}
此外，该同构与满射的复合 $I\surj J\surj K$ 相容。

特别地，取 $J$ 为单点集，便得到
\begin{equation*}
\Hk^{I}_{G,\bK}|_{\D(U)}\cong \Hk^{\{1\}}_{G,\bK}.
\end{equation*}

同样的性质对 $\Sht^{I}_{G,\bK}$ 也成立。若 $\th: I\surj J$ 是有限集之间的满射，则存在典范同构
\begin{equation*}
\Sht^{I}_{G,\bK}|_{\D_{\th,U}(U^{J})}\cong \Sht^{J}_{G,\bK}.
\end{equation*}
此外，该同构与满射的复合 $I\surj J\surj K$ 相容。

特别地，取 $J$ 为单点集，便得到
\begin{equation*}
\Sht^{I}_{G,\bK}|_{\D(U)}\cong \Sht^{\{1\}}_{G,\bK}.
\end{equation*}

\subsection{水平结构的改变}\label{ss:change level}

\ssszh{水平提升}

对每个 $x\in \Sig$，设 $\bK^{\#}_{x}\lhd \bK_{x}$ 是正规子群，其商 $\bL_{x}=\bK_{x}/\bK^{\#}_{x}$ 有限。令 $\bK^{\#}=(\bK^{\#}_{x})_{x\in \Sig}$。则自然映射
\begin{equation*}
\Sht^{(I_{1},\cdots, I_{r})}_{G,\bK^{\#}} \to \Sht^{(I_{1},\cdots, I_{r})}_{G,\bK}
\end{equation*}
是 $L$-主丛，其中 $L:=\prod_{x\in \Sig}\bL_{x}(k_{x})$ 为有限群。

对这样的正规子群 $\bK^{\#}$ 取投影极限，便得到 $U^{I}$ 上的叠
\begin{equation*}
\Sht^{(I_{1},\cdots, I_{r})}_{G,\infty\Sig} :=\varprojlim_{\bK^{\#}}\Sht^{(I_{1},\cdots, I_{r})}_{G,\bK^{\#}}.
\end{equation*}
其投影
\begin{equation*}
\Sht^{(I_{1},\cdots, I_{r})}_{G,\infty\Sig}\to \Sht^{(I_{1},\cdots, I_{r})}_{G,\bK}
\end{equation*}
是 pro-\'etale 映射，且是以 $K=\prod_{x\in \Sig}K_{x}$ 为结构群的主丛，其中 $K_{x}=\bK_{x}(k_{x})$。特别地有：

\begin{lem}\label{l:change level Sig}
对每个 $x\in \Sig$，设 $\bK'_{x}\subset \bK_{x}$ 是水平群，并令 $\bK'=(\bK'_{x})_{x\in \Sig}$。则自然映射
\begin{equation*}
\Sht^{(I_{1},\cdots, I_{r})}_{G,\bK'} \to \Sht^{(I_{1},\cdots, I_{r})}_{G,\bK}
\end{equation*}
是有限 \'etale 的，其几何纤维与陪集 $\prod_{x\in \Sig}K_{x}/K'_{x}$ 一一对应（其中 $K'_{x}=\bK'_{x}(k_{x})$）。
\end{lem}

\ssszh{在 $x$ 处的满水平结构}\label{sss:full x}

若 $x\notin \Sig$，则可定义 $(U-\{x\})^{I}$ 上的叠
\begin{equation*}
\Sht^{(I_{1},\cdots, I_{r})}_{G,\bK,\infty x}=\varprojlim_{\bK^{\da}_{x}}\Sht^{(I_{1},\cdots, I_{r})}_{G,\bK^{\da}}
\end{equation*}
其中 $\bK^{\da}$ 表示沿 $\{x\}\cup \Sig$ 的水平结构，即在 $\Sig$ 上取 $\bK$、在 $x$ 处取 $\bK^{\da}_{x}$。这是 $\Sht^{(I_{1},\cdots, I_{r})}_{G, \bK}|_{(U-\{x\})^{I}}$ 上的 $G(\cO_{x})$-主丛。

我们断言，$G(\cO_{x})$ 在 $\Sht^{(I_{1},\cdots, I_{r})}_{G, \bK,\infty x}$ 上的作用可典范地延拓为 $G(F_{x})$ 的作用。事实上，设
\begin{equation*}
\xi=((x_{i})_{i\in I}, (\cE_{j})_{0\le j\le r}, (\a_{j})_{1\le j\le r}, \io, (\ka_{j})_{0\le j\le r})
\end{equation*}
是 $\Sht^{(I_{1},\cdots, I_{r})}_{G,\bK,\infty x}$ 的一个 $R$-点，其中 $\ka_{j}: G\times D_{x, R}\isom \cE_{j}|_{D_{x,R}}$ 是对应于 $x$ 处满水平结构的平凡化。这里 $D_{x,R}=\Spec(\cO_{x}\wh\ot R)$，$D_{x,R}^{\times}=\Spec(F_{x}\wh\ot R)$。设 $g\in G(F_{x})$，定义 $\xi\cdot g$ 为多元组
\begin{equation*}
((x_{i})_{i\in I}, (\cE^{g}_{j})_{0\le j\le r}, (\a'_{j})_{1\le j\le r}, \io', (\ka'_{j})_{0\le j\le r}).
\end{equation*}
其中 $\cE^{g}_{j}$ 是把 $\cE_{j}|_{(X-\{x\})_{R}}$ 与 $D_{x,R}$ 上的平凡 $G$-主丛用同构
\begin{equation*}
G\times D^{\times}_{x, R}\xr{g} G\times D^{\times}_{x, R} \xr{\ka_{j}} \cE_{j}|_{D^{\times}_{x,R}}.
\end{equation*}
粘合所得的 $G$-主丛。

上式第一个映射是 $g$ 的左乘。映射 $\a'_{j}$ 与 $\io'$ 由 $\a_{j}$ 与 $\io$ 诱导得到，因为按定义 $\G(x_{j})$ 与 $\{x\}_{R}$ 不交。最后，$\cE'_{j}|_{D_{x,R}}$ 的平凡化 $\ka'_{j}$ 就是由 $\cE'_{j}$ 的粘合构造给出的重言平凡化。

\subsection{中心作用}\label{sss:central action}

设 $Z$ 是 $G$ 的中心。对每个 $x\in \Sig$，令 $\bK_{Z,x}=\bK_{x}\cap L^{+}_{x}Z$，这是 $Z$ 在 $x$ 处的水平群。记 $\bK_{Z}$ 为诸 $\bK_{Z,x}$（$x\in \Sig$）的汇集。由于 $Z$ 交换，叠 $\Bun_{Z,\bK_{Z}}$ 是群叠。特别地，$\Bun_{Z,\bK_{Z}}(k)$ 是 Picard 群胚。

$\Bun_{Z,\bK_{Z}}$ 在 $\Bun_{G,\bK}$ 上有典范作用。为看出这一点，先考虑如下情形：对每个 $x\in \Sig$，$\bK_{x}$ 是主同余子群 $L^{(n_{x})}_{x}G$，其中 $n_{x}\in\NN$。此时 $\Bun_{G,\bK}=\Bun_{G,D}$，其中 $D=\sum n_{x}x$，且 $\Bun_{Z,\bK_{Z}}=\Bun_{Z,D}$。

考虑 $S$-点 $\cE\in\Bun_{G,D}(S)$，它对应于 $X\times S$ 上的右 $G$-主丛 $\cE$ 连同 $\cE|_{D\times S}$ 的平凡化 $\io_{D}$。再考虑 $S$-点 $\cQ\in \Bun_{Z,D}(S)$，它对应于 $X\times S$ 上的右 $Z$-主丛 $\cQ$ 连同 $\cQ|_{D\times S}$ 的平凡化 $\ka_{D}$。由于 $Z$ 在 $G$ 中中心，它在 $\cE$ 上的自然作用与 $G$ 的作用交换，故 $Z$ 通过 $G$-主丛自同构作用于 $\cE$。

作扭纤维积 ${}_{\cQ}\cE:=\cQ\times^{Z}_{X\times S}\cE$，即 $\cQ\times_{X\times S}\cE$ 对 $Z$ 的对角作用的商。$\cE$ 上的右 $G$-作用使 ${}_{\cQ}\cE$ 具有 $G$-主丛结构。利用 $\cE|_{D\times S}$ 与 $\cQ|_{D\times S}$ 的平凡化 $\io_{D}$ 与 $\ka_{D}$，得到 ${}_{\cQ}\cE|_{D\times S}$ 的平凡化 $\io'_{D}$。于是定义作用 $(\cQ,\ka_{D})\cdot (\cE,\io_{D})$ 的结果为 $({}_{\cQ}\cE, \io'_{D})$。

对更一般的水平结构，仍可如下得到所需的作用。对每个 $x$ 选取足够小的主同余子群 $L^{(n_{x})}_{x}G\subset \bK_{x}$，并如 \S\ref{sss:level} 那样把 $\Bun_{G,\bK}$（相应地 $\Bun_{Z,\bK_{Z}}$）实现为 $\Bun_{G,D}$（相应地 $\Bun_{Z,D}$）对 $G_{D}$（相应地 $Z_{D}$）的某个代数子群的商。细节从略。

现在考虑上述构造的 shtuka 版本。我们断言，$\Bun_{Z,\bK_{Z}}(k)$ 在 $\Sht^{(I_{1},\cdots, I_{r})}_{G,\bK}$ 上有典范作用。事实上，对任意 $S$-点
\begin{equation*}
\xi=((x_{i})_{i\in I}, (\cE_{j})_{0\le j\le r}, (\a_{j})_{1\le j\le r}, \io)
\end{equation*}
（属于 $\Sht^{(I_{1},\cdots, I_{r})}_{G,\bK}$）以及 $\cQ\in \Bun_{Z,\bK_{Z}}(k)$，用上述构造得到 ${}_{\cQ}\cE_{j}\in \Bun_{G,\bK}(S)$。

原有的同构 $\a_{j}: \cE_{j-1}|_{X\times S-\cup_{i\in I_{j}}\G(x_{i})}\isom  \cE_{j}|_{X\times S-\cup_{i\in I_{j}}\G(x_{i})}$ 诱导同构
\begin{equation*}
\a'_{j}: {}_{\cQ}\cE_{j-1}|_{X\times S-\cup_{i\in I_{j}}\G(x_{i})}\isom  {}_{\cQ}\cE_{j}|_{X\times S-\cup_{i\in I_{j}}\G(x_{i})},
\end{equation*}
而 $\io:\cE_{r}\isom{}^{\t}\cE_{0}$ 诱导同构
\begin{equation*}
\io': {}_{\cQ}\cE_{r}\cong {}_{{}^{\t}\cQ}({}^{\t}\cE_{0})={}^{\t}({}_{\cQ}\cE_{0}),
\end{equation*}
因为 ${}^{\t}\cQ=\cQ$。于是 $\cQ$ 在 $\Sht^{(I_{1},\cdots, I_{r})}_{G,\bK}$ 上的作用把 $\xi$ 送到
\begin{equation*}
\xi':=((x_{i})_{i\in I}, ({}_{\cQ}\cE_{j})_{0\le j\le r}, (\a'_{j})_{1\le j\le r}, \io')\text{。}
\end{equation*}

容易看出，对任意支配余权组 $\un\l\in (\xcoch(T)^{+})^{I}$，$\Bun_{Z,\bK_{Z}}(k)$ 在 $\Sht^{(I_{1},\cdots, I_{r})}_{G,\bK}$ 上的作用保持有界版本 $\Sht^{(I_{1},\cdots, I_{r}),\le\un\l}_{G,\bK}$。

\begin{exam} 考虑 $G=\GL_{n}$ 且无水平结构的情形。此时 $\Bun_{Z}=\Pic_{X}$，而线丛 $\cL\in \Pic_{X}(k)$ 在 $\Sht^{(I_{1},\cdots, I_{r})}_{\GL_{n}}$ 上的作用就是把每个向量丛 $\cV_{j}$ 与 $\cL$ 作张量积。

另一方面，考虑 $G=\SL_{n}$ 且无水平结构的情形。此时 $\Bun_{Z}\cong \Pic_{X}[n]$，即 $X$ 的 Picard 叠的 $n$-挠子叠。对 $\cL\in \Bun_{Z}(k)=\Pic_{X}[n](k)$，即带同构 $\e: \cO_{X}\cong \cL^{\ot n}$ 的线丛 $\cL$，它在 $\Sht^{(I_{1},\cdots, I_{r})}_{\SL_{n}}$ 上的作用仍由把每个向量丛 $\cV_{j}$ 与 $\cL$ 作张量积给出。注意，$\det(\cV_{j})$ 原有的平凡化连同 $\e$ 诱导了 $\det(\cV_{j}\ot \cL)\cong \det(\cV_{j})\ot \cL^{\ot n}$ 的平凡化。
\end{exam}

\subsection{Hecke 对应}

\ssszh{球面 Hecke 对应}

设 $x\notin \Sig$，$K_{x}=G(\cO_{x})$，并取 $[g_{x}]\in K_{x}\bs G(F_{x})/K_{x}$。定义
\begin{equation*}
 \Sht^{(I_{1},\cdots, I_{r})}_{G,\bK}[g_{x}]:= \Sht^{(I_{1},\cdots, I_{r})}_{G,\bK,\infty x}/(K_{x}\cap g_{x}K_{x}g_{x}^{-1}).
\end{equation*}
我们把 $\Sht^{(I_{1},\cdots, I_{r})}_{G,\bK}[g_{x}]$ 看作对应
\begin{equation*}
\begin{tikzcd}
& \Sht^{(I_{1},\cdots, I_{r})}_{G,\bK}[g_{x}] \arrow[dl, "\oll{c}"'] \arrow[dr, "\orr{c}"] & \\
\Sht^{(I_{1},\cdots, I_{r})}_{G,\bK}|_{(U-\{x\})^{I}} & & \Sht^{(I_{1},\cdots, I_{r})}_{G,\bK}|_{(U-\{x\})^{I}}
\end{tikzcd}
\end{equation*}
其中 $\oll{c}$ 是对应于包含 $K_{x}\cap g_{x}K_{x}g_{x}^{-1}$ 的自然投影；映射 $\orr{c}$ 是复合
\begin{equation*}
 \Sht^{(I_{1},\cdots, I_{r})}_{G,\bK,\infty x}/(K_{x}\cap g_{x}K_{x}g_{x}^{-1})\xr{\cdot g_{x}} \Sht^{(I_{1},\cdots, I_{r})}_{G,\bK,\infty x}/(g_{x}^{-1}K_{x}g_{x}\cap K_{x})\to \Sht^{(I_{1},\cdots, I_{r})}_{G,\bK}.
\end{equation*}
这里第一个映射用到 \S\ref{sss:full x} 中构造的右 $G(F_{x})$-作用；第二个映射对应于包含 $g^{-1}_{x}K_{x}g_{x}\cap K_{x}\subset K_{x}$。

\ssszh{分歧点处的 Hecke 对应}

对 $x\in \Sig$ 与 $[g_{x}]\in  K_{x}\bs G(F_{x})/K_{x}$，定义对应
\begin{equation}\label{Hk x Sig}
\begin{tikzcd}
& \Sht^{(I_{1},\cdots, I_{r})}_{G,\bK}[g_{x}] \arrow[dl, "\oll{c}"'] \arrow[dr, "\orr{c}"] & \\
\Sht^{(I_{1},\cdots, I_{r})}_{G,\bK} & &  \Sht^{(I_{1},\cdots, I_{r})}_{G,\bK}
\end{tikzcd}
\end{equation}
为此令
\begin{equation*}
\Sht^{(I_{1},\cdots, I_{r})}_{G,\bK}[g_{x}]:=\Sht^{(I_{1},\cdots, I_{r})}_{G,\bK^{x}, \infty x}/(K_{x}\cap g_{x}K_{x}g_{x}^{-1}).
\end{equation*}
这里 $\bK^{x}$ 是水平群组 $(\bK_{x})_{x\in \Sig-\{x\}}$。映射 $\oll{c}$ 是自然投影，而映射 $\orr{c}$ 是先作用 $g_{x}$、再取自然投影，与球面 Hecke 修正的情形相同。

对 $k$-概形 $S$，$\Sht^{(I_{1},\cdots, I_{r})}_{G,\bK}[g_{x}](S)$ 分类 $(x_{i})_{i\in I}$（其中 $x_{i}\in U(S)$）连同图表
\begin{equation*}
\begin{tikzcd}[column sep=large, row sep=large]
\cE_{0} \arrow[d, dashed, "f_{0}"] \arrow[r, dashed, "\a_{1}"] &
\cE_{1} \arrow[d, dashed, "f_{1}"] \arrow[r, dashed, "\a_{2}"] &
\cdots \arrow[r, dashed, "\a_{r}"] &
\cE_{r} \arrow[d, dashed, "f_{r}"] \arrow[r, "\io", "\sim"'] &
{}^{\t}\cE_{0} \arrow[d, dashed, "{}^{\t}f_{0}"] \\
\cE'_{0} \arrow[r, dashed, "\a'_{1}"] &
\cE'_{1} \arrow[r, dashed, "\a'_{1}"] &
\cdots \arrow[r, dashed, "\a'_{r}"] &
\cE'_{r} \arrow[r, "\io'", "\sim"'] &
{}^{\t}\cE'_{0}
\end{tikzcd}
\end{equation*}
这里上下两行都是 $\Sht^{(I_{1},\cdots, I_{r})}_{G,\bK}$ 的 $S$-点，支点为 $(x_{i})_{i\in I}$；竖直箭头是 $U\times S$ 上 $G$-主丛的同构
\begin{equation*}
f_{i}: \cE_{i}|_{U\times S}\isom \cE'_{i}|_{U\times S},
\end{equation*}
它可延拓为 $x'\in \Sigma-\{x\}$ 处 $\bK_{x'}$-水平结构的同构，并在 $x$ 处有相对位置 $[g_{x}]\in K_{x}\bs G(F_{x})/K_{x}$。

\ssszh{有界版本}

给定界 $\un\l\in (\xcoch(T)^{+})^{I}$，有带界 $\un\l$ 的 Hecke 对应的有界版本
\begin{equation*}
\Sht^{(I_{1},\cdots, I_{r}), \le \un\l}_{G,\bK}[g_{x}]
\end{equation*}
它视为 $\Sht^{(I_{1},\cdots, I_{r}), \le \un\l}_{G,\bK}|_{(U-\{x\})^{I}}$（若 $x\notin \Sig$）或 $\Sht^{(I_{1},\cdots, I_{r}), \le \un\l}_{G,\bK}$（若 $x\in \Sig$）的自对应。

\subsection{部分 Frobenius}

存在典范态射
\begin{equation}\label{pl Fr}
\Fr^{\pl}: \Sht^{(I_{1},\cdots, I_{r})}_{G,\bK}\to \Sht^{(I_{2},\cdots, I_{r}, I_{1})}_{G,\bK}
\end{equation}
它在 $S$-点层面的定义如下：把 $S$-点
\begin{equation*}
\begin{tikzcd}[column sep=small]
\xi=((x_{i})_{i\in I}, \quad \cE_{0} \arrow[r, dashed, "\a_{1}"] & \cdots \arrow[r, dashed, "\a_{r-1}"] & \cE_{r-1} \arrow[r, dashed, "\a_{r}"] & \cE_{r} \arrow[r, "\io", "\sim"'] & {}^{\t}\cE_{0}
\end{tikzcd}
\end{equation*}
送到
\begin{equation*}
\begin{tikzcd}[column sep=small]
\Fr^{\pl}(\xi)=( (x'_{i})_{i\in I}, \quad \cE_{1} \arrow[r, dashed, "\a_{2}"] & \cdots \arrow[r, dashed, "\a_{r}"] & \cE_{r} \arrow[r, dashed, "{{}^{\t}\a_{1}\c\io}"] & {}^{\t}\cE_{1} \arrow[r, equal, "{\io'=\id}"] & {}^{\t}\cE_{1}.
\end{tikzcd}
\end{equation*}
这里
\begin{equation*}
x'_{i}=\begin{cases} x_{i}\c\Fr_{S}, & i\in I_{1}\\
x_{i}, & i\notin I_{1}.\end{cases}
\end{equation*}
显然
\begin{equation*}
\underbrace{\Fr^{\pl}\c \Fr^{\pl}\c\cdots \c \Fr^{\pl}}_{r \text{ 次}}= \Fr\in \End(\Sht^{(I_{1},\cdots, I_{r})}_{G,\bK}).
\end{equation*}

有交换图
\begin{equation}\label{pl Fr UI}
\begin{tikzcd}
\Sht^{(I_{1},\cdots, I_{r})}_{G,\bK} \arrow[r, "\Fr^{\pl}"] \arrow[d, "\pi^{(I_{1},\cdots, I_{r})}"] & \Sht^{(I_{2},\cdots, I_{1})}_{G,\bK} \arrow[d, "\pi^{(I_{2},\cdots, I_{1})}"] \\
U^{I} \arrow[r, "\Fr_{I_{1}}"] & U^{I}
\end{tikzcd}
\end{equation}
这里 $\Fr_{I_{1}}:U^{I}\to U^{I}$ 是在 $I_{1}$ 的各分量上取 Frobenius、在其余分量上取恒等的映射。由于 $U^{I}$ 上的 $\Fr_{I_{1}}$ 与 $\Fr^{\pl}$ 都是万有同胚，该图是拓扑笛卡尔图。

\section{上同调上的结构}\label{s:coho}

本节总结 $G$-shtuka 模叠的上同调（相对于支点空间而言）上的基本对称性：分解性、Hecke 代数作用与部分 Frobenius 作用。在 C.Xue 的讲义 \cite{Xue} 中，这些对称性将用于进一步研究 $G$-shtuka 模叠的上同调。

\subsection{Shtuka 模叠的上同调}

\ssszh{Satake 层}

回忆 Lusztig、Beilinson--Drinfeld、Ginzburg、Mirkovi\'c--Vilonen 的几何 Satake 等价（见 \cite{MV}）。有 $\Qlbar$ 上的张量范畴等价
\begin{equation*}
\Sat: \Rep(\dG,\Qlbar)\isom \Perv_{\Aut(D)}(\QGr_{G})
\end{equation*}
其中右端的幺半结构由卷积给出，交换约束由 Beilinson 与 Drinfeld 的融合构造给出。这里 $\QGr_{G}$ 上的 $\Aut(D)$-等变反常层，指的是 $\Gr_{G}$ 上 $L^{+}G\rtimes\Aut(D)$-等变的反常层。对任意有限集 $I$，定义函子
\begin{equation*}
\Sat_{I}: \Rep(\dG^{I},\Qlbar)\to\Perv((\QGr_{G}/\Aut(D))^{I}).
\end{equation*}
对 $V=\boxtimes_{i\in I}V_{i}$（其中 $V_{i}\in \Rep(\dG)$），定义 $\Sat_{I}(V)=\boxtimes_{i\in I}\Sat(V_{i})$，再按可加性延拓。

下面在 $\Hk^{I}_{G,\bK}$ 上定义层的复形。设 $r=|I|$，并取双射 $\s: I\isom \{1,2,\cdots, r\}$。令 $I_{j}=\{\s^{-1}(j)\}\subset I$，便得到把 $I$ 分成单点的分解 $I=I_{1}\coprod\cdots\coprod I_{r}$。回忆 \eqref{ev Hk singleton} 中的相对位置映射
\begin{equation*}
\ev^{(I_{1},\cdots, I_{r})}_{\Hk}: \Hk^{(I_{1},\cdots, I_{r})}_{G,\bK}\to (\QGr_{G}/\Aut(D))^{I}.
\end{equation*}
定义拉回
\begin{equation*}
\IC^{(I_{1},\cdots, I_{r})}_{\Hk}(V):=(\ev^{(I_{1},\cdots, I_{r})}_{\Hk})^{*}\Sat_{I}(V)\in D_{c}(\Hk^{(I_{1},\cdots, I_{r})}_{G,\bK}, \Qlbar).
\end{equation*}
若 $V=\boxtimes_{i\in I}V_{i}$ 且每个 $V_{i}$ 不可约、最高权为 $\l_{i}\in\xcoch(T)^{+}$，则 $\IC^{(I_{1},\cdots, I_{r})}_{\Hk}(V)$ 支撑在闭子叠 $\Hk^{(I_{1},\cdots, I_{r}), \le\un\l}_{G,\bK}$ 上，其中 $\un\l=(\l_{i})_{i\in I}$。

又有投影
\begin{equation*}
\pi^{(I_{1},\cdots, I_{r})}_{I}: \Hk^{(I_{1},\cdots, I_{r})}_{G,\bK}\to \Hk^{I}_{G,\bK}.
\end{equation*}
于是定义
\begin{equation*}
\IC^{I,\s}_{\Hk}(V):=\pi^{(I_{1},\cdots, I_{r})}_{I, !}\IC^{(I_{1},\cdots, I_{r})}_{\Hk}(V)\in D_{c}(\Hk^{I}_{G,\bK}, \Qlbar).
\end{equation*}
同样，若 $V=\boxtimes_{i\in I}V_{i}$ 且 $V_{i}$ 不可约、最高权为 $\l_{i}\in\xcoch(T)^{+}$，则 $\IC^{I}_{\Hk}(V)$ 支撑在 $\Hk^{I, \le\un\l}_{G,\bK}$ 上，其中 $\un\l=(\l_{i})_{i\in I}$。

\begin{lem}\label{l:IC Hk}
对任意 $V\in \Rep(\dG^{I})$，复形 $\IC^{I,\s}_{\Hk}(V)$ 在相差一个典范同构的意义下与 $I$ 上顺序 $\s$ 的选取无关。我们记之为 $\IC^{I}_{\Hk}(V)$。

此外，若 $V=\boxtimes_{i\in I}V_{i}$ 且 $V_{i}$ 不可约、最高权为 $\l_{i}$，则 $\IC^{I,\s}_{\Hk}(V)$ 在相差一个平移的意义下同构于 $\Hk^{I, \le\un\l}_{G,\bK}$ 的相交复形，其中 $\un\l=(\l_{i})_{i\in I}$。
\end{lem}
\begin{proof}[证明概要] 设 $V=\boxtimes_{i\in I}V_{i}$ 且 $V_{i}$ 不可约、最高权为 $\l_{i}$。对分成单点的分解 $I=I_{1}\coprod\cdots\coprod I_{r}$，由命题 \ref{p:fm sm}，相对位置映射 $\ev^{(I_{1},\cdots, I_{r})}_{\Hk}$ 形式光滑。因此 $\IC^{(I_{1},\cdots, I_{r})}_{\Hk}(V)$ 在相差一个平移的意义下典范同构于 $\Hk^{(I_{1},\cdots, I_{r}), \le\un\l}_{G,\bK}$ 的相交复形。映射 $\pi^{(I_{1},\cdots, I_{r})}_{I}$ 在 \cite[命题 4.2 之后]{MV} 的意义下分层小，这意味着 $\IC^{I,\s}_{\Hk}(V)$ 在相差一个平移的意义下典范同构于 $\Hk^{I, \le\un\l}_{G,\bK}$ 的相交复形。故它与 $\s$ 的选取无关。
\end{proof}

更一般地，若 $I=I_{1}\coprod\cdots \coprod I_{r}$ 是 $I$ 的任意分解，则可给每个 $I_{j}$ 排序，并把这分解细化为分成单点的分解 $I=I'_{1}\coprod\cdots \coprod I'_{s}$。于是有双射 $\s: I\isom \{1,2,\cdots, s\}$，使 $I'_{j}=\{\s^{-1}(j)\}$（$1\le j\le s$），且 $\s(I_{j})$ 是 $\{1,2,\cdots, s\}$ 的连续子集。有典范投影
\begin{equation*}
\pi^{(I'_{1},\cdots, I'_{s})}_{(I_{1},\cdots, I_{r})}: \Hk^{(I'_{1},\cdots, I'_{s})}_{G,\bK}\to \Hk^{(I_{1},\cdots, I_{r})}_{G,\bK}.
\end{equation*}
对 $V\in \Rep(\dG^{I}, \Qlbar)$ 定义
\begin{equation*}
\IC^{(I_{1},\cdots, I_{r}),\s}_{\Hk}(V):=\pi^{(I'_{1},\cdots, I'_{s})}_{(I_{1},\cdots, I_{r}), !}\IC^{(I'_{1},\cdots, I'_{s})}_{\Hk}(V)\in D_{c}(\Hk^{(I_{1},\cdots, I_{r})}_{G,\bK}, \Qlbar).
\end{equation*}

与引理 \ref{l:IC Hk} 相同的论证给出以下结论。

\begin{lem}\label{l:IC Hk refine} 设 $V\in \Rep(\dG^{I})$。
\begin{enumerate}
\item $\Hk^{(I_{1},\cdots, I_{r})}_{G,\bK}$ 上的复形 $\IC^{(I_{1},\cdots, I_{r}),\s}_{\Hk}(V)$ 在相差一个典范同构的意义下，与细化分解 $I=I_{1}\coprod\cdots \coprod I_{r}$ 的 $I$ 上顺序 $\s$ 的选取无关。记此复形为 $\IC^{(I_{1},\cdots, I_{r})}_{\Hk}(V)$。

\item 设 $I=I'_{1}\coprod\cdots \coprod I'_{r'}$ 是分解 $I=I_{1}\coprod\cdots \coprod I_{r}$ 的任意细化。则存在典范同构
\begin{equation*}
\IC^{(I_{1},\cdots, I_{r})}_{\Hk}(V)\cong \pi^{(I'_{1},\cdots, I'_{r'})}_{(I_{1},\cdots, I_{r}), !}\IC^{(I'_{1},\cdots, I'_{r'})}_{\Hk}(V).
\end{equation*}
且这些同构与进一步的细化相容。

\item 若 $V=\boxtimes_{i}V_{i}$ 且每个 $V_{i}$ 不可约、最高权为 $\l_{i}$，则 $\IC^{(I_{1},\cdots, I_{r})}_{\Hk}(V)$ 在相差一个平移的意义下同构于 $\Hk^{(I_{1},\cdots, I_{r}),\le\un\l}_{G,\bK}$ 的相交复形，其中 $\un\l=(\l_{i})_{i\in I}$。
\end{enumerate}
\end{lem}

\ssszh{$\Sht^{I}_{G}$ 的上同调复形}

回忆典范映射
\begin{equation*}
\r^{(I_{1},\cdots, I_{r})}: \Sht^{(I_{1},\cdots, I_{r})}_{G,\bK}\to \Hk^{(I_{1},\cdots, I_{r})}_{G,\bK}.
\end{equation*}
对 $V\in \Rep(\dG^{I}, \Qlbar)$ 定义
\begin{equation*}
\IC^{(I_{1},\cdots, I_{r})}_{\Sht}(V):=(\r^{(I_{1},\cdots, I_{r})})^{*}\IC^{(I_{1},\cdots, I_{r})}_{\Hk}(V)\in D_{c}(\Sht^{(I_{1},\cdots, I_{r})}_{G,\bK}).
\end{equation*}

另一种做法是：当 $I=I_{1}\coprod\cdots \coprod I_{r}$ 是分成单点的分解时，可考虑 $\Sht^{(I_{1},\cdots, I_{r})}_{G,\bK}$ 的相对位置映射 \eqref{ev Sht singleton}。于是
\begin{equation*}
\IC^{(I_{1},\cdots, I_{r})}_{\Sht}(V)\cong (\ev_{\Sht}^{(I_{1},\cdots, I_{r})})^{*}\Sat_{I}(V).
\end{equation*}

\begin{lem}\label{l:IC Sht refine} 设 $V\in \Rep(\dG^{I})$。
\begin{enumerate}
\item 对分解 $I=I_{1}\coprod\cdots \coprod I_{r}$ 的任意细化 $I=I'_{1}\coprod\cdots \coprod I'_{r'}$，存在典范同构
\begin{equation*}
\IC^{(I_{1},\cdots, I_{r})}_{\Sht}(V)\cong \pi^{(I'_{1},\cdots, I'_{r'})}_{(I_{1},\cdots, I_{r}), !}\IC^{(I'_{1},\cdots, I'_{r'})}_{\Sht}(V),
\end{equation*}
且与进一步的细化相容。
\item 若 $V=\boxtimes_{i\in I}V_{i}$，其中 $V_{i}$ 不可约、最高权为 $\l_{i}$，则 $\IC^{(I_{1},\cdots, I_{r})}_{\Sht}(V)$ 在相差 $\#I$ 的平移的意义下同构于 $\Sht^{(I_{1},\cdots, I_{r}), \le\un\l}_{G,\bK}$ 的相交复形，其中 $\un\l=(\l_{i})_{i\in I}$。
\end{enumerate}
\end{lem}
\begin{proof}
(1) 由引理 \ref{l:IC Hk refine} 以及笛卡尔方图
\begin{equation*}
\begin{tikzcd}[column sep=large, row sep=large]
\Sht^{(I'_{1},\cdots, I'_{r'})}_{G,\bK} \arrow[r, "\r^{(I'_{1},\cdots, I'_{r'})}"] \arrow[d, "\pi^{(I'_{1},\cdots, I'_{r'})}_{(I_{1},\cdots, I_{r})}"] & \Hk^{(I'_{1},\cdots, I'_{r'})}_{G,\bK} \arrow[d, "\pi^{(I'_{1},\cdots, I'_{r'})}_{(I_{1},\cdots, I_{r})}"] \\
\Sht^{(I_{1},\cdots, I_{r})}_{G,\bK} \arrow[r, "\r^{(I_{1},\cdots, I_{r})}"] & \Hk^{(I_{1},\cdots, I_{r})}_{G,\bK}
\end{tikzcd}
\end{equation*}
的固有基变换得到。

(2) 先设各 $I_{j}$ 都是单点集，此时结论由命题 \ref{p:fm sm} 得到。

现考虑一般的分解 $I=I_{1}\coprod\cdots\coprod I_{r}$。把每个 $I_{j}$ 分成单点，再合起来便得到分解 $I=I'_{1}\coprod\cdots\coprod I'_{r'}$，其中每个 $I'_{i}$ 是单点，且 $I_{j}$ 是若干连续部分的并。由命题 \ref{p:local model}，映射 $\pi: \Sht^{(I'_{1},\cdots, I'_{r'}),\le\un\l}_{G,\bK}\to\Sht^{(I_{1},\cdots, I_{r}),\le\un\l}_{G,\bK}$ 在 \'etale 局部同构于映射 $\Gr^{(I'_{1},\cdots, I'_{r'}),\un\le\l}_{G}\to \Gr^{(I_{1},\cdots, I_{r}),\un\le\l}_{G}$，而后者分层小。故由 (1)，
\begin{equation*}
\begin{aligned}
&\IC^{(I_{1},\cdots, I_{r})}_{\Sht}(V)\cong \pi_{!}\IC^{(I'_{1},\cdots, I'_{r'})}_{\Sht}(V)\\
&\qquad\cong\pi_{!}\IC(\Sht^{(I'_{1},\cdots, I'_{r'}),\le\un\l}_{G,\bK})[-\#I]
\cong\IC(\Sht^{(I_{1},\cdots, I_{r}),\le\un\l}_{G,\bK})[-\#I].
\end{aligned}
\end{equation*}
\end{proof}

\begin{dfn} 设 $I$ 是带分解 $I=I_{1}\coprod\cdots \coprod I_{r}$ 的有限集。回忆映射 $\pi^{I}: \Sht^{(I_{1},\cdots, I_{r})}_{G,\bK}\to U^{I}$。定义 ind-可构造复形
\begin{equation*}
\cH^{(I_{1},\cdots, I_{r})}_{G,\bK}(V):=\pi^{I}_{!}\IC^{(I_{1},\cdots, I_{r})}_{\Sht}(V)\in D(U^{I}, \Qlbar).
\end{equation*}
当讨论中 $(G,\bK)$ 固定时，就把 $\cH^{(I_{1},\cdots, I_{r})}_{G,\bK}(V)$ 简记为 $\cH^{(I_{1},\cdots, I_{r})}(V)$。

特别地，有 ind-可构造复形 $\cH^{I}(V)\in D(U^{I}, \Qlbar)$。
\end{dfn}

\begin{rmk}
为看出 $\cH^{I}(V)$ 是 ind-可构造的，我们使用 \S\ref{sss:HN} 中讨论的 Harder--Narasimhan 截断。一般水平结构可归结为主同余水平结构的情形，故不妨设 $\bK_{x}\subset L^{+}_{x}G$ 对一切 $x\in \Sig$ 成立。对 $\mu\in \HN$，令 ${}^{\le \mu}\pi^{I}: {}^{\le\mu}\Sht^{I}_{G,\bK}\to U^{I}$ 为 $\pi^{I}$ 的限制。令 ${}^{\le \mu}\IC_{\Sht}^{I}(V)$ 为 $\IC^{I}_{\Sht}(V)$ 到 ${}^{\le\mu}\Sht^{I}_{G,\bK}$ 的限制。则 ${}^{\le\mu}\cH^{I}(V):={}^{\le \mu}\pi^{I}_{!}({}^{\le \mu}\IC_{\Sht}^{I}(V))$ 可构造，因为 ${}^{\le \mu}\IC_{\Sht}^{I}(V)$ 的支撑在 $U^{I}$ 上有限型。于是有
\begin{equation}\label{HV colim}
\cH^{I}(V)=\colim_{\mu\in \HN}{}^{\le\mu}\cH^{I}(V)
\end{equation}
作为 $U^{I}$ 上的 ind-可构造复形。这里对 $\mu\le \mu'$，过渡映射 ${}^{\le\mu}\cH^{I}(V)\to {}^{\le\mu'}\cH^{I}(V)$ 由开浸入 ${}^{\le\mu}\Sht^{I}_{G,\bK}\incl {}^{\le\mu'}\Sht^{I}_{G,\bK}$ 诱导。
\end{rmk}

\begin{cor}[引理 \ref{l:IC Sht refine} 的推论]\label{c:same HV}
对分解 $I=I_{1}\coprod\cdots \coprod I_{r}$ 的任意细化 $I=I'_{1}\coprod\cdots \coprod I'_{r'}$，存在典范同构
\begin{equation*}
\cH^{(I_{1},\cdots, I_{r})}(V)\cong \cH^{(I'_{1},\cdots, I'_{r'})}(V),
\end{equation*}
且与进一步的细化相容。特别地，对任意分解 $I=I_{1}\coprod\cdots \coprod I_{r}$，存在典范同构
\begin{equation*}
\cH^{(I_{1},\cdots ,I_{r})}(V)\cong\cH^{I}(V)\in D(U^{I}, \Qlbar).
\end{equation*}
\end{cor}
\begin{proof}
对引理 \ref{l:IC Sht refine}(1) 中的同构作用 $\pi^{(I_{1},\cdots, I_{r})}_{!}$。
\end{proof}
